\chapter{Algebraic preliminaries}\label{ch:prelim}
This chapter should be thought as a companion for the geometric applications in Chapter \ref{ch:ffuk} and \ref{ch:approx}. The definitions and results appearing in this chapter are mostly due to other authors. We decided to introduce the theory of $dg$-categories and persistence categories through the lenses of enriched category theory, which the author learned in \cite{Kelly2005Basic, Fong_Spivak_2019}. \S\ref{se:dgcat} is a rewriting of the basic standard theory of $dg$-categories (that the author learned in \cite{KellerICM,Se:book-fukaya-categ}) while \S\ref{se:catpm} and \S\ref{se:pc} contains an introduction to the theory of persistence categories (originally introduced in \cite{BCZ:tpc}).


\S\ref{se:Ainf} contains an introduction to filtered $A_\infty$-categories. Most of the material presented there is by now standard, with the exception of the discussion of of $dg$-functors (appearing in \S\ref{dgpullpush}), filtered $A_\infty$-functors (appearing in \S\ref{ap:pshf}) and persistence Hochschild homology (appearing in \S\ref{app:PHH}). These results are due to the author and collaborators and were originally established in \cite{ABC26}.

\S\ref{se:TPC} contains a quick introduction to the theory of triangulated persistence categories (TPCs), as originally introduced in \cite{BCZ:tpc}. In \S\ref{se:approx} we define the concept of \say{categorical metric approximability} for objects of TPCs. This concept was introduced in \cite{ABC26} and is a quantitative generalization of the concept of generation and split-generation in triangulated categories, and can be seen - in a sense to be explained - as a genralization of the concept of total boundedness in the theory of metric spaces.

Finally, in \S\ref{se:sys} we describe the notion of a \emph{system of categories with increasing accuracy}, as defined in \cite{ABC26}. Although technically involved, this framework is currently necessary for the study of persistence Fukaya categories.

\section{Introduction}
$A_\infty$-categories were introduced by Fukaya in \cite{Fu:morse-homotopy-ai} as a homotopical generalization of $dg$-categories. In turn, $dg$-categories arose as a response to foundational difficulties in the theory of triangulated categories (see, for example, the survey \cite{KellerICM}). Triangulated categories, originally introduced by Verdier to formalize derived categories, appear naturally in many areas of mathematics. However, it is well known that they exhibit several undesirable categorical properties. 

For instance, the bicategory of triangulated categories does not admit a monoidal structure. This deficiency reflects a more fundamental issue: mapping cones are not functorial. Concretely, given a triangulated category $\Tau$, there is no cone functor from the arrow category of $\Tau$ to $\Tau$.

The introduction of $dg$-categories was motivated by the desire to remedy these shortcomings. While this approach is successful from several perspectives (again, see \cite{KellerICM}), it comes at the cost of increased complexity. $dg$-categories encode a substantial amount of information, and $A_\infty$-categories even more so. Consequently, when working with a (pre-triangulated) $A_\infty$-category, one often ultimately discards part of the structure and passes to the associated triangulated homotopy category.

Filtered $dg$-categories and filtered $A_\infty$-categories are natural refinements of their unfiltered counterparts, in which the algebraic structures are required to be compatible with a filtration on the morphism complexes indexed by the real line. In symplectic topology, the motivation for such structures is clear: Floer complexes are naturally filtered chain complexes whenever the action functional is defined (see Chapter~\ref{ch:ffuk}). 

Just as the homology of a filtered chain complex naturally carries the structure of a persistence module, the homotopy category of a filtered $dg$- or $A_\infty$-category inherits the structure of a \say{persistence category}, namely a category enriched over the category of persistence modules.